% mage-auto-heading: what-a-stack-is
% mage-auto-heading: the-seven-and-how-they-relate
% mage-auto-heading: where-to-look
% mage-auto-heading: why-stacks-exist
% Authoritative LaTeX manuscript; Markdown is a generated compatibility projection.

A single mechanism can eliminate one failure class. A useful engineering capability often requires several mechanisms.

A provenance record tells us who changed an artifact, but not whether every mutation was recorded. A model can make an engineering fact explicit, but nothing guarantees that consumers read the current version or that the implementation still matches it. A sensor detects a failure; detection alone does not restore service. So mechanisms compose.

\section{What a Stack Is}\label{heading:what-a-stack-is}

An \textbf{engineering stack is a recurring set of mechanisms that together provide a useful engineering capability.} The stacks here are reference architectures, not prescriptions: a different substrate may provide the same capability with fewer, different, or stronger mechanisms.

Start with the capability:

% <!-- inline-quote -->

\begin{quote}
What must the environment be able to do, and which guarantees are required to make that claim true?
\end{quote}

The seven stacks that follow answer that question for recurring engineering concerns.

\section{The Seven, and How They Relate}\label{heading:the-seven-and-how-they-relate}

\begin{figure}[htbp]
\centering
\includegraphics[width=\linewidth]{book/latex/assets/appendix-a-capability-map-1a4bbe8ef7.pdf}
\caption{Seven reference engineering stacks and their common relationships. Arrows show relationships among capabilities that recurred in one system, not a required architecture or adoption order. Resource Mediation and Context Delivery are shown separately because they provide useful capabilities independently.}
\label{appendix-a-capability-map}
\end{figure}

\section{Where to Look}\label{heading:where-to-look}

A mechanism addresses a specific obligation or failure. A move names the transferable engineering judgment behind related mechanisms. A stack combines mechanisms into a reusable capability. Together, the stacks contribute to a governed engineering environment.

This appendix starts from a capability and asks which mechanisms must work together to provide it. Appendix B starts from recurring problems and shows how the same engineering move can produce different mechanisms in different settings. Each stack gives its capability, constituent moves, dependencies, and required composition. Concrete implementations of these moves appear in the companion MAGE Mechanism Catalog.

\section{Why Stacks Exist}\label{heading:why-stacks-exist}

A stack is one form engineering capital can take when several mechanisms are needed to provide a capability. One failure may motivate a sensor. Another may expose a bypass that warrants a gate. A third may reveal missing provenance that warrants attribution. Mechanisms introduced separately can eventually depend on one another strongly enough to form a coherent stack.
